\section{Men of Destiny}

Remembering oneself, discovering the true will, following one's Dharma, can all be summarized by the ancient call to \emph{know thyself}. This is what defies every attempt to define or limit Tradition. It is not a question of reading the next book, knowing about Etruscan fertility goddesses, or adopting the very next craze. It is irrelevant to a debate about which religion is traditional or not. It is not about adhering to any particular political program. Tradition is not reducible to a formula.

There are very specific spiritual states: detachment, a calm mind, mental strength, or a feeling of deep-rooted security. That is the only standard and the necessary pre-condition to understand Tradition in any meaningful sense. If you are anxious, if your mind is scattered, if you react emotionally, then you need to follow some spiritual practice that will lead to such the state of awareness he describes. Becoming a pagan, or Buddhist, or Hindu, or thelemite is insufficient in itself, and those who convert often reveal a disheveled, superficial mind. \textbf{Julius Evola} explains:

\begin{quotationx}
Yet that is not enough. It is not by chance that the formula, ``know thyself'', which, in its deepest meaning, refers exactly to such teachings, was itself written on the Delphic temple of Apollo, i.e., the Hyperborean God\footnote{\url{https://gornahoor.net/Hyperborean.htm}}. Letting such traditional truths act on oneself, until they awaken precise interior forces, means to proceed along the way that leads to a spiritual level at which the meaning of life constitutes something absolutely different from what the rest of men think: a sense of clarity, of absolute strength, of incomparable security.

But to have an inkling of all that, to glimpse a ``character'' in whom a type of indomitability is united to the feeling of detachment of ``those who are joined from afar'' and to an interior inaccessibility; in whom, therefore, there is simultaneously a superior calm, an aloofness, and a readiness to attack, to command, to absolute action---to have demonstrated this ``character'', means also to have foreshadowed the mystery of the primordial Nordic race, or Hyperborean race, as the race of the spirit. Such is in fact the olympian and solar mode of being; the popular imagination refers to them today as the so-called ``men of destiny'' and yesterday referred to them as the vanished types of great leaders---in reality, in that it is now about the last echoes or flashing of what was characteristic, in general, of the great Hyperborean super race, before its dispersion and alteration. One remembers Plutarch's expression about the members of the ancient Roman Senate: ``They sit like a council of kings.''

Hence, a further consequence: if a civilization of the ``classical'' type, in the olympian and virile sense, not in the vulgar aesthetic and formulaic sense, reflects something of the Nordic race of the spirit, every romantic and ``tragic'' civilization, as opposed to it, will instead be the certain mark of the prevalence of the earlier influence by races and ethnic residues of a non-Nordic, pre-Aryan and anti-Aryan nature.
\end{quotationx}


Despite all the talk about the Law of Cycles, it is poorly understood. It means that the material and historical circumstances have deteriorated to a certain level. By the Law of Elective Affinity, certain spirits will be attracted to that world and will feel very natural in it. Coincidentally, as the Wheel turns ever so slightly\footnote{\url{https://gornahoor.net/images/circlespirito.jpg}} to begin the next cycle, there will be attracted a smaller group. Their task is to be the spiritual elite for the next cycle, the men of destiny. To do so, these men will have to make special efforts to remember, to recall the distant past, to truly know themselves. Those who recall their Past will create the Future of Tradition.


\flrightit{Posted on 2012-04-12 by Cologero}

\begin{center}* * *\end{center}

\begin{footnotesize}
\begin{sffamily}
\texttt{Cologero on 2012-04-17 at 17:37 said:}

The few readers who take this seriously may be interested in comparing this passage from \textit{The Consolations of Philosophy}. Here, Sophia is speaking to Boethius:

\begin{quotex}\footnotesize

Now I know the other cause, or rather the major cause of your illness: \textbf{you have forgotten your true nature}. And so I have found out in full the reason for your sickness and the way to approach the task of restoring you to health. It is because you are confused by loss of memory that you wept and claimed you you had been banished and robbed of all your possessions. And it is because you don't know the end and purpose of things that you think the wicked and the criminal have power and happiness. And because you have forgotten the means by which the world is governed you believe these ups and downs of fortune happen haphazardly. These are grave causes and they lead not only to illness but even death. Thanks, however, to the Author of all health, nature had not quite abandoned you. In your true belief about the world government--that it is subject to divine reason and not the haphazards of change--there lies our greatest hope of rekindling your health.
\end{quotex}


\end{sffamily}
\end{footnotesize}

